\documentclass[12pt]{article}
\usepackage[utf8]{inputenc}
\usepackage[T1]{fontenc}
\usepackage{hyperref}
\title{Eblaite Language Resources for K-12 Curriculum Development}
\author{Compiled by Claude Sonnet 5.5 (claude-sonnet-5-5) for Andr\'as Kornai}
\date{}
\begin{document}
\maketitle

\section{Introduction}
Eblaite (ISO 639-3 xeb; Glottocode ebla1238) is an extinct Semitic language of the third millennium BC in northern Syria. It is known only from the cuneiform tablets of the Ebla palace archive. Glottolog classifies it as East Semitic. Scholars disagree on whether it is a separate language or a variety of Old Akkadian, and on where it sits within Semitic. There are no native speakers. The usable corpus is fairly small.

\section{Primary Educational Resources}
I found no primers, textbooks or school programmes for Eblaite. Any K-12 heritage use would have to be built from scholarly sources.

\section{Linguistic and Reference Materials}
\begin{itemize}
\item Ebla tablets (Wikipedia): \url{https://en.wikipedia.org/wiki/Ebla_tablets}
\item Eblaite language (Wikipedia): \url{https://en.wikipedia.org/wiki/Eblaite_language}
\item Eblaite language (Britannica): \url{https://www.britannica.com/print/article/177618}
\item Pettinato, \emph{Old Canaanite Cuneiform Texts of the Third Millennium} (Undena, 1979), described in search results as the first grammar of Eblaite: \url{https://undena.com/EL-UP/Pettinato_1979_Old_Canaanite_Cuneiform_Texts_-_MANE_1.7.pdf}
\item Article on the Ebla chancery texts in light of recent linguistic research (summary page): \url{https://plosjournal.deepdyve.com/lp/de-gruyter/the-ebla-chancery-texts-in-the-light-of-recent-linguistic-research-th0Gdl3DRP}
\item Glottolog entry: \url{https://glottolog.org/resource/languoid/id/ebla1238.bigmap.html}
\end{itemize}

\section{Religious and Community Resources}
None found in the search results.

\section{Digital and Audio Resources}
\begin{itemize}
\item CDLI record for an Ebla tablet (P244156): \url{https://cdli.earth/artifacts/244156}
\item The Palace Archive of Ebla, Syria: \url{https://historyofinformation.com/detail.php?id=3426}
\end{itemize}
I found no audio resources.

\section{Curriculum Development Context}
Eblaite is a historical-heritage topic. Work with it would center on the writing system, the archive and the Sumerian--Eblaite bilingual word lists. It would not aim at speaking the language, and the scholarly debates should be presented honestly.

\end{document}